% GENERATED FILE. Edit canonical lessons, then run npm run generate:books.
% canonical-source-hash: fnv1a64:f77a19f06923db44
% canonical-lessons: GU-C178-saru-karvu, GU-C178-patavvu, GU-C178-samaysar, GU-C178-baki, GU-C178-vitvu

\chapter{Starting and finishing}
\label{ch:gu-starting-and-finishing}
\hlchaptermodality{178}

\begin{chapteropening}
\textbf{By the end of this chapter:} \emph{I can say when I will start and finish something.}

Complete the last lesson of chapter 178: I can say when I will start and finish something.
\end{chapteropening}

\section[\texorpdfstring{\gu{શરૂ} \gu{કરવું} (śarū karvũ)}{શરૂ કરવું (śarū karvũ)}]{\gu{શરૂ} \gu{કરવું} (śarū karvũ) — to start, to begin}

\label{lesson:GU-C178-saru-karvu}



\begin{quote}
Before the new one: say the Gujarati for regular, then the Gujarati for everyday.
\end{quote}

\subsection*{You'll want to know: \gu{શરૂ} \gu{કરવું}}
\textbf{\gu{શરૂ} \gu{કરવું}} — \emph{śarū karvũ} — \textquotedblleft{}to start, to begin\textquotedblright{}.

\begin{quote}
\textbf{\gu{હું} \gu{કાલે} \gu{નવું} \gu{કામ} \gu{શરૂ} \gu{કરીશ}.} — \emph{hũ kāle navũ kām śarū karīś} — \textquotedblleft{}I will start a new job tomorrow.\textquotedblright{}
\end{quote}

\begin{cousinweb}
\textbf{\gu{શરૂ} \gu{કરવું}} (\emph{śarū karvũ}) is built from \emph{śarū}, \textquotedblleft{}begun\textquotedblright{}, and the inherited \emph{karvũ}, \textquotedblleft{}to do\textquotedblright{}. \emph{Śarū} is a loan from Arabic \emph{shurūʿ}, \textquotedblleft{}a beginning\textquotedblright{}, through Persian, the word behind \emph{śarūāt} (Lane).
\end{cousinweb}

\begin{grammarlens}[title={What you've built}]
One more word for everyday life at A2.
\end{grammarlens}

\subsection*{Your turn}
\begin{itemize}
\raggedright
  \item \emph{Say it:} \emph{śarū karvũ}
  \item \emph{Say it:} \emph{śarū karvũ}, once more
  \item \emph{Say it:} \emph{hũ kāle navũ kām śarū karīś} — I will start a new job tomorrow.
  \item \emph{Recall:} say the Gujarati for regular, then the Gujarati for everyday, then say \emph{śarū karvũ} again
\end{itemize}

\begin{tcolorbox}[breakable,colback=teal!4,colframe=teal!35!black,title={Before you move on}]
What does \gu{શરૂ} \gu{કરવું} mean? (\textquotedblleft{}To start, to begin\textquotedblright{}.) Say it once more.
\end{tcolorbox}
\section[\texorpdfstring{\gu{પતાવવું} (patāvvũ)}{પતાવવું (patāvvũ)}]{\gu{પતાવવું} (patāvvũ) — to finish, to settle}

\label{lesson:GU-C178-patavvu}



\begin{quote}
Before the new one: say the Gujarati for everyday, then the Gujarati for to start.
\end{quote}

\subsection*{You'll want to know: \gu{પતાવવું}}
\textbf{\gu{પતાવવું}} — \emph{patāvvũ} — \textquotedblleft{}to finish, to settle\textquotedblright{}.

\begin{quote}
\textbf{\gu{તે} \gu{આજે} \gu{કામ} \gu{પતાવશે}.} — \emph{te āje kām patāvśe} — \textquotedblleft{}He will finish the work today.\textquotedblright{}
\end{quote}

\begin{cousinweb}
\textbf{\gu{પતાવવું}} (\emph{patāvvũ}) is a tadbhava. CDIAL lists it, glossed \textquotedblleft{}to have done with, settle\textquotedblright{}, under Sanskrit \emph{prāpta}, \textquotedblleft{}reached, obtained\textquotedblright{} (CDIAL 8945).
\end{cousinweb}

\begin{grammarlens}[title={What you've built}]
One more word for everyday life at A2.
\end{grammarlens}

\subsection*{Your turn}
\begin{itemize}
\raggedright
  \item \emph{Say it:} \emph{patāvvũ}
  \item \emph{Say it:} \emph{patāvvũ}, once more
  \item \emph{Say it:} \emph{te āje kām patāvśe} — He will finish the work today.
  \item \emph{Recall:} say the Gujarati for everyday, then the Gujarati for to start, then say \emph{patāvvũ} again
\end{itemize}

\begin{tcolorbox}[breakable,colback=teal!4,colframe=teal!35!black,title={Before you move on}]
What does \gu{પતાવવું} mean? (\textquotedblleft{}To finish, to settle\textquotedblright{}.) Say it once more.
\end{tcolorbox}
\section[\texorpdfstring{\gu{સમયસર} (samaysar)}{સમયસર (samaysar)}]{\gu{સમયસર} (samaysar) — on time}

\label{lesson:GU-C178-samaysar}



\begin{quote}
Before the new one: say the Gujarati for to start, then the Gujarati for to finish.
\end{quote}

\subsection*{You'll want to know: \gu{સમયસર}}
\textbf{\gu{સમયસર}} — \emph{samaysar} — \textquotedblleft{}on time\textquotedblright{}.

\begin{quote}
\textbf{\gu{હું} \gu{સમયસર} \gu{આવીશ}.} — \emph{hũ samaysar āvīś} — \textquotedblleft{}I will come on time.\textquotedblright{}
\end{quote}

\begin{cousinweb}
\textbf{\gu{સમયસર}} (\emph{samaysar}) is built from the taught \emph{samay}, \textquotedblleft{}time\textquotedblright{}, a tatsama from Sanskrit \emph{samaya}, and the ending \emph{-sar}, \textquotedblleft{}in keeping with\textquotedblright{}.
\end{cousinweb}

\begin{grammarlens}[title={What you've built}]
One more word for everyday life at A2.
\end{grammarlens}

\subsection*{Your turn}
\begin{itemize}
\raggedright
  \item \emph{Say it:} \emph{samaysar}
  \item \emph{Say it:} \emph{samaysar}, once more
  \item \emph{Say it:} \emph{hũ samaysar āvīś} — I will come on time.
  \item \emph{Recall:} say the Gujarati for to start, then the Gujarati for to finish, then say \emph{samaysar} again
\end{itemize}

\begin{tcolorbox}[breakable,colback=teal!4,colframe=teal!35!black,title={Before you move on}]
What does \gu{સમયસર} mean? (\textquotedblleft{}On time\textquotedblright{}.) Say it once more.
\end{tcolorbox}
\section[\texorpdfstring{\gu{બાકી} (bākī)}{બાકી (bākī)}]{\gu{બાકી} (bākī) — remaining, left over}

\label{lesson:GU-C178-baki}



\begin{quote}
Before the new one: say the Gujarati for to finish, then the Gujarati for on time.
\end{quote}

\subsection*{You'll want to know: \gu{બાકી}}
\textbf{\gu{બાકી}} — \emph{bākī} — \textquotedblleft{}remaining, left over\textquotedblright{}.

\begin{quote}
\textbf{\gu{હજી} \gu{કામ} \gu{બાકી} \gu{છે}.} — \emph{hajī kām bākī che} — \textquotedblleft{}There is still work left.\textquotedblright{}
\end{quote}

\begin{cousinweb}
\textbf{\gu{બાકી}} (\emph{bākī}) is a loan from Arabic \emph{bāqī}, \textquotedblleft{}remaining, lasting\textquotedblright{}, which reached Gujarati through Persian. Its Arabic root is \emph{b-q-y}, \textquotedblleft{}to remain\textquotedblright{} (Lane).
\end{cousinweb}

\begin{grammarlens}[title={What you've built}]
One more word for everyday life at A2.
\end{grammarlens}

\subsection*{Your turn}
\begin{itemize}
\raggedright
  \item \emph{Say it:} \emph{bākī}
  \item \emph{Say it:} \emph{bākī}, once more
  \item \emph{Say it:} \emph{hajī kām bākī che} — There is still work left.
  \item \emph{Recall:} say the Gujarati for to finish, then the Gujarati for on time, then say \emph{bākī} again
\end{itemize}

\begin{tcolorbox}[breakable,colback=teal!4,colframe=teal!35!black,title={Before you move on}]
What does \gu{બાકી} mean? (\textquotedblleft{}Remaining, left over\textquotedblright{}.) Say it once more.
\end{tcolorbox}
\section[\texorpdfstring{\gu{વીતવું} (vītvũ)}{વીતવું (vītvũ)}]{\gu{વીતવું} (vītvũ) — to go by, to pass (of time)}

\label{lesson:GU-C178-vitvu}



\begin{quote}
Before the new one: say the Gujarati for on time, then the Gujarati for remaining.
\end{quote}

\subsection*{You'll want to know: \gu{વીતવું}}
\textbf{\gu{વીતવું}} — \emph{vītvũ} — \textquotedblleft{}to go by, to pass (of time)\textquotedblright{}.

\begin{quote}
\textbf{\gu{સમય} \gu{જલદી} \gu{વીતશે}.} — \emph{samay jaldī vītśe} — \textquotedblleft{}Time will pass quickly.\textquotedblright{}
\end{quote}

\begin{cousinweb}
\textbf{\gu{વીતવું}} (\emph{vītvũ}) is a tadbhava. It comes from Sanskrit \emph{vṛtta}, \textquotedblleft{}turned; passed, elapsed (of time)\textquotedblright{}, and CDIAL glosses the Gujarati verb \textquotedblleft{}to pass (of time)\textquotedblright{} (CDIAL 12069).
\end{cousinweb}

\begin{grammarlens}[title={What you've built}]
One more word for everyday life at A2.
\end{grammarlens}

\subsection*{Your turn}
\begin{itemize}
\raggedright
  \item \emph{Say it:} \emph{vītvũ}
  \item \emph{Say it:} \emph{vītvũ}, once more
  \item \emph{Say it:} \emph{samay jaldī vītśe} — Time will pass quickly.
  \item \emph{Recall:} say the Gujarati for on time, then the Gujarati for remaining, then say \emph{vītvũ} again
\end{itemize}

\begin{tcolorbox}[breakable,colback=teal!4,colframe=teal!35!black,title={Before you move on}]
What does \gu{વીતવું} mean? (\textquotedblleft{}To go by, to pass (of time)\textquotedblright{}.) Say it once more.
\end{tcolorbox}
